% ============================================================
%  ABSTRACT
% ============================================================
\begin{abstract}

\draft{Deploying legged mobile manipulators for autonomous emergency
assessment requires the robot to actively adapt its body posture to the
perceptual demands of each mission phase---locating, approaching, scanning,
and performing ultrasound on an unknown patient.}

This paper presents \ac{TERESA}, an autonomous emergency assessment system on the
Boston Dynamics Spot quadruped with a Unitree Z1 arm (\draft{hereafter Spotzi}). \draft{\ac{TERESA} treats body posture as an active perceptual
resource, adapting it to the demands of four sequential phases:
(1) the robot searches for the patient by walking forward while
sweeping body posture; (2) it approaches while concurrently scanning
the torso, eliminating idle time; (3) it executes an exposure scan
covering the body surface and detecting injuries; and (4) it performs
\ac{FAST} ultrasound with pre-planned posture optimisation for each
anatomical target. A web-based interface allows the \draft{remote} operator to click
any scanned body point and trigger a full-body reconfiguration to
revisit it.}

\draft{Experiments on the Spotzi platform in a laboratory
environment, with a medical training manikin, demonstrate fully autonomous
execution of all four phases.}
\end{abstract}
